% !TEX program = xelatex
% Argus 商业计划书（天使/种子轮） — 编译： xelatex Argus_BP.tex （跑两遍生成目录与页码）
\documentclass[UTF8,11pt,fontset=none]{ctexart}

%==================== 字体 ====================
\usepackage{fontspec}
\setCJKmainfont{Noto Serif CJK SC}
\setCJKsansfont{Noto Sans CJK SC}
\setCJKmonofont{Noto Sans Mono CJK SC}
\setmainfont{Latin Modern Roman}
\setsansfont{Latin Modern Sans}

%==================== 版式 ====================
\usepackage[a4paper,margin=2.3cm,top=2.4cm,bottom=2.4cm]{geometry}
\usepackage{xcolor}
\usepackage{booktabs}
\usepackage{tabularx}
\usepackage{array}
\usepackage{colortbl}
\usepackage{enumitem}
\usepackage{titlesec}
\usepackage{tcolorbox}
\usepackage{fontawesome5}
\usepackage{graphicx}
\usepackage{ragged2e}
\usepackage[hidelinks]{hyperref}

\tcbuselibrary{skins,breakable}

%==================== 品牌色 ====================
\definecolor{argusnavy}{HTML}{14304F}
\definecolor{argusteal}{HTML}{0E7C86}
\definecolor{argusaccent}{HTML}{E76F51}
\definecolor{argusgray}{HTML}{4A5568}
\definecolor{arguslight}{HTML}{EAF1F4}
\definecolor{argusline}{HTML}{C7D3DB}

\hypersetup{linkcolor=argusteal,urlcolor=argusteal,citecolor=argusteal}

%==================== 标题样式 ====================
\titleformat{\section}
  {\sffamily\Large\bfseries\color{argusnavy}}
  {\colorbox{argusnavy}{\color{white}\;\thesection\;}}{0.6em}{}
\titleformat{\subsection}
  {\sffamily\large\bfseries\color{argusteal}}
  {\thesubsection}{0.6em}{}
\titlespacing*{\section}{0pt}{1.4em}{0.7em}
\titlespacing*{\subsection}{0pt}{1.0em}{0.4em}

\setlist[itemize]{leftmargin=1.4em,itemsep=2pt,topsep=2pt}
\setlist[enumerate]{leftmargin=1.6em,itemsep=2pt,topsep=2pt}
\renewcommand{\arraystretch}{1.25}

% 段首不缩进、段间留白（商务文档风格）
\setlength{\parindent}{0pt}
\setlength{\parskip}{0.5em}

%==================== 自定义盒子 ====================
\newtcolorbox{keybox}[1][]{
  enhanced, breakable, colback=arguslight, colframe=argusteal,
  left=10pt, right=10pt, top=6pt, bottom=6pt,
  arc=2pt, boxrule=0pt, leftrule=3pt, #1}

\newtcolorbox{quotebox}[1][]{
  enhanced, colback=white, colframe=argusaccent,
  left=10pt, right=10pt, top=6pt, bottom=6pt,
  arc=1pt, boxrule=0pt, leftrule=4pt, #1}

\newcommand{\tagline}[1]{{\sffamily\color{argusgray}#1}}
\newcommand{\hl}[1]{\textbf{\color{argusnavy}#1}}

% 表头风格
\newcommand{\thd}[1]{\textbf{\color{white}#1}}
\newcommand{\headrow}{\rowcolor{argusnavy}}
\newcommand{\altrow}{\rowcolor{arguslight}}

\begin{document}
\sffamily

%======================================================================
% 封面
%======================================================================
\begin{titlepage}
\thispagestyle{empty}
\vspace*{1.2cm}
\begin{flushleft}
{\color{argusline}\rule{\textwidth}{2pt}}\\[1.0cm]
{\sffamily\fontsize{56}{60}\selectfont\bfseries\color{argusnavy} ARGUS}\\[0.3cm]
{\sffamily\Large\color{argusteal} 自主 AI 科学家 \;/\; Autonomous AI Scientist}\\[0.8cm]
{\color{argusline}\rule{\textwidth}{1pt}}\\[1.2cm]

{\sffamily\huge\bfseries\color{argusnavy} 商业计划书}\\[0.5cm]
{\sffamily\Large\color{argusgray} 天使 / 种子轮融资}\\[2.0cm]

\begin{quotebox}
{\large 给它一个研究目标和机器规则，剩下的它自己来：\\
选题 → 调研 → 设计实验 → 在\textbf{真实 benchmark} 上跑 → 分析 → 写论文 → 自审 → 改稿 → 投稿就绪。\\[4pt]
\textbf{没有人在回路里逐步审批。}}
\end{quotebox}

\vfill
{\sffamily\color{argusgray}
日期：2026 年 6 月 \hfill 文件密级：\textbf{机密 · Confidential}\\
版本：v1.0 \hfill 联系方式：\_\_\_\_\_\_\_\_\_\_\_\_\_\_\_\_}\\[0.4cm]
{\color{argusline}\rule{\textwidth}{2pt}}
\end{flushleft}
\end{titlepage}

%======================================================================
% 目录
%======================================================================
\thispagestyle{empty}
{\sffamily\color{argusnavy}\Large\bfseries 目录}\\[0.3em]
{\color{argusline}\rule{\textwidth}{1pt}}
\vspace{0.4em}
\setcounter{tocdepth}{1}
\tableofcontents
\clearpage

%======================================================================
% 执行摘要
%======================================================================
\section{执行摘要}

\hl{Argus 是一个 7×24 全自主运行的"AI 科学家"。} 用户只需提供两样东西——一个\textbf{研究目标}和一份\textbf{机器规则}（GPU、路径、预算等约束）——Argus 就独立地把一篇论文从无到有做出来：选题、文献调研、设计实验、在真实公开 benchmark 上完整跑、分析结果、撰写 LaTeX 论文、自我评审、反复改稿，直到达到投稿就绪。\textbf{科研所需的每一个判断都由 AI agent 自己下，没有人在回路里逐步审批。}

\begin{keybox}
\textbf{一句话定位：} 别的研究 agent 只能做到论文的"前 70\%"（一个点子 + 一个玩具实验 + 一份看上去像样的 PDF）。Argus 专攻评审真正会拒稿的"后 30\%"——证据不足、实验单薄、结论无据、引用造假、版面超标、产物失效、不可复现——靠\textbf{持续执行}加\textbf{硬性质量闸门}把这 30\% 补上。
\end{keybox}

\textbf{为什么是现在（Why now）：} 前沿大模型已经能写代码、能用工具、能长时间运行，但单个模型的上下文窗口装不下一个跨越数百步、数亿 token 的长程科研任务。\textbf{真正的瓶颈不再是"模型够不够聪明"，而是"如何让一个会遗忘的智能体，在超长任务里保持连贯、不造假、不空转"。} 这正是 Argus 的工程内核所解决的问题。

\textbf{我们已经建成了什么（事实，非愿景）：}
\begin{itemize}
  \item 完整的\textbf{三层 Agent 架构}（Planner / Engineer / Reviewer）+ 一根领域无关的"笨管道"harness，已可 7×24 daemon 化无人值守运行。
  \item \textbf{8 阶段研究流水线}（research → plan → benchmark → run → analysis → draft → review → submission），每阶段由 Reviewer 对照 checklist 裁决是否推进。
  \item \textbf{65 个内置 Skill}（55 个 Engineer + 10 个 Reviewer），覆盖文献检索、实验执行、消融、结论溯源、论文写作、图表生成、格式预检、投稿保障。
  \item \textbf{防造假诚信护栏}：强制使用真实公开 benchmark、禁止灌水、留存可审计产物——这是产品的核心，不是附加功能。
  \item \textbf{长程记忆 / 反遗忘上下文引擎}：跨 session 只交接"经筛选的有价值记忆"，根治长任务里的 amnesia 空转——这是我们最深的技术壁垒。
  \item 工程成熟度：153 个测试文件、CI（pytest/ruff/mypy）全绿、支持多会议格式（EMNLP/ACL、AAAI）。
\end{itemize}

\textbf{商业机会：} 全球每年研发支出超过 2.5 万亿美元，每年产出数百万篇学术论文；在企业研发、量化金融、生物医药、AI 实验室中，"高水平研究人力"是最稀缺、最昂贵的资源。Argus 把"一篇可投稿论文"的边际成本从\textbf{数月人力（数万美元）}降到\textbf{数百美元的算力}，是 10–100 倍的成本结构重构。

\textbf{融资诉求：} 本轮拟融资 \hl{人民币 1,500 万元}（约 US\$2.1M），出让 15\%–18\%，用于研发团队扩张、算力/API、商业化落地与合规。18 个月内冲刺首篇端到端认证论文、10 家种子客户、ARR 突破人民币 300 万元，达到 A 轮就绪。

%======================================================================
% 问题
%======================================================================
\section{问题：科研生产力的结构性瓶颈}

\textbf{研究是人类最有价值、却最难规模化的劳动。} 一篇高水平论文背后是数月的文献阅读、实验设计、调参跑数、结果分析、反复写作与改稿。这条链路上的每一环都依赖稀缺的高水平人力，且高度串行、难以并行复制。

\textbf{现有的"AI 辅助科研"工具只解决了边角问题：}
\begin{itemize}
  \item \hl{文献类工具}（Elicit、Consensus、Scite 等）只做检索与综述，不碰实验与论文成稿。
  \item \hl{通用 Coding Agent}（Devin、通用 Codex/Claude 调用等）能写代码，但不理解科研全流程，没有证据闸门，跑着跑着就开始编造结果、或在长任务里"失忆"反复空转。
  \item \hl{学术界的"AI Scientist"原型}能演示端到端，但停在"前 70\%"：玩具实验、薄弱证据、看上去像样但经不起评审的 PDF。
\end{itemize}

\begin{quotebox}
\textbf{核心痛点：} 把一个聪明的大模型，变成一个\textbf{能在数百步、数亿 token 的长程任务里保持连贯、不造假、不空转，并产出可被同行评审接受的成果}的系统——这件事至今没有人真正做成。难点不在"生成"，而在\textbf{长程执行的工程治理}。
\end{quotebox}

\textbf{三个被长期低估的技术难点：}
\begin{enumerate}
  \item \hl{长程记忆 / 上下文遗忘。} 单模型上下文窗口有限；长任务里 session 被反复有损压缩上百次，模型丢失工作记忆后会反复重读同一批资料空转（amnesia loop）。
  \item \hl{诚信与可复现。} 没有硬约束时，agent 倾向于"造一个好看的结果"而非"做一个真实的实验"——伪造图表、灌水数据行、引用占位作者。
  \item \hl{完成度判定。} "这篇研究到底做完没有 / 够不够格"无法用关键词或正则判断，需要一个真正会读产物、会下判断的评审者。
\end{enumerate}

%======================================================================
% 解决方案
%======================================================================
\section{解决方案：Argus —— 自主 AI 科学家}

Argus 的设计建立在一句反直觉但关键的原则上：

\begin{keybox}
\textbf{"管道（harness）没有 agent 自己聪明。"}\\[4pt]
所以我们把系统严格切成两层：一切\textbf{科研判断}（这个 idea 平不平庸？证据够不够？该投了吗？）都交给 \textbf{Agent}；harness 只做\textbf{领域无关的脏活}（预算、调度、持久化、结构化 I/O、防造假护栏）。每一次"用词面去二次猜 agent"的诱惑，我们都删掉，换成结构化信号或交还给 agent——因为只要 harness 开始替 agent 做判断，它就成了系统能力的天花板。
\end{keybox}

这条哲学是 Argus 区别于其他"流水线式"研究 agent 的根本：\textbf{Argus 不是把 prompt 串起来的脚本，而是一个能自己做判断的自主研究工厂。} 它的直接收益是——随着底层大模型变强，Argus 的能力上限自动抬高，而不被 harness 里写死的规则卡住。

%======================================================================
% 产品与技术
%======================================================================
\section{产品与技术}

\subsection{三层 Agent + 笨管道}

\begin{tabularx}{\textwidth}{@{}l l X@{}}
\toprule
\headrow \thd{层} & \thd{角色} & \thd{职责} \\
\midrule
\altrow L4 & Planner & backlog 空了排下一个任务；项目认证不达标时自动改派认证任务。 \\
L1 & Engineer & 单轮执行：搜文献、写代码、跑实验、写 LaTeX；与 Reviewer 共享工作目录。 \\
\altrow L2 & Reviewer & 对照当前阶段 checklist 审查，给 \texttt{done}/\texttt{continue}（附具体修复指令）/\texttt{blocked}。\textbf{是项目是否完成的唯一事实来源。} \\
\bottomrule
\end{tabularx}

三个 agent 都基于前沿推理模型（当前默认 OpenAI \texttt{gpt-5.4}），Reviewer 拥有 shell 访问权，能自己读文件、跑检查，而非依赖 harness 喂结论。harness 负责预算/限速、持久化与记忆、调度与 daemon、结构化 I/O 和防造假护栏——\textbf{只做这些，且只做这些}。

\subsection{8 阶段研究流水线}

\begin{center}
\colorbox{argusteal}{\color{white}\sffamily\;research\;}$\rightarrow$
\colorbox{argusteal}{\color{white}\sffamily\;plan\;}$\rightarrow$
\colorbox{argusteal}{\color{white}\sffamily\;benchmark\;}$\rightarrow$
\colorbox{argusteal}{\color{white}\sffamily\;run\;}$\rightarrow$
\colorbox{argusteal}{\color{white}\sffamily\;analysis\;}$\rightarrow$
\colorbox{argusteal}{\color{white}\sffamily\;draft\;}$\rightarrow$
\colorbox{argusteal}{\color{white}\sffamily\;review\;}$\rightarrow$
\colorbox{argusnavy}{\color{white}\sffamily\;submission\;}
\end{center}

每个阶段的产物由 Engineer 产出，由 Reviewer 对照该阶段 checklist 裁决是否推进。\textbf{关键在于"硬性证据闸门"}——例如：必须有全部方法/基线的完整打分行才能出论文；每一个标题数字结论都必须能溯源到原始产物；每一张图都要有真实生成的栅格图、prompt/输出哈希与审查记录；引用必须查真实性，拦截占位作者与断链。这些闸门正是把竞品的"前 70\%"推到"可投稿"的关键。

\subsection{长程记忆 / 反遗忘上下文引擎（核心壁垒）}

这是 Argus 最深、最难被复制的工程资产，也直接回应了第 2 节的头号技术难点。

\begin{keybox}
长任务的 Codex session 会涨到数亿 token、被有损压缩上百次，每次压缩都丢工作记忆，导致模型反复重读同一批文档\textbf{空转（amnesia loop）}。Argus 的修复哲学不是"加看门狗"，而是让 \textbf{session 结构上短命 + 跨 session 边界只交接"经筛选的有价值记忆"}。
\end{keybox}

具体机制：
\begin{itemize}
  \item \hl{受控工作记忆检查点（Curated Working Memory）}：一份\textbf{有硬上限}的小状态（目标 / 已完成 / 试过但失败 / 当前阻塞 / 下一步）。上限在代码层强制——\textbf{"遗忘"本身就是一种主动筛选}，地面真相留在磁盘产物里可随时召回。
  \item \hl{Reviewer 即记忆审计员}：Engineer 在回合末尾提出记忆更新提案，由 Reviewer 对照实际产物校验后落定为下一份权威检查点（提议—验证分离，防止模型自欺）。
  \item \hl{结构化 Session 滚动}：一个会话活满 N 轮就主动退役、开新会话，只把检查点带过去——从结构上杜绝长 session 的退化。
\end{itemize}

\textbf{这套机制可迁移到一切"长程自主 agent"场景}（不止科研）：长周期软件工程、运维、金融研究、法律尽调——任何"任务远超上下文窗口"的领域，都需要这层记忆治理。它是 Argus 从"科研工具"走向"长程 agent 平台"的技术底座。

\subsection{诚信护栏：防造假是产品核心}

唯一正当的"硬规则"是\textbf{反作弊}：必须用真实公开 benchmark、禁止重复行灌水、强制留存可审计包。这些约束的是"作弊"而非"科研选择"，所以留在 harness 里合理。对企业客户而言，\textbf{"可审计、可复现、不造假"是把 AI 研究产物用于真实决策的前提}——这是 Argus 相对学术原型的关键商业护城河。

\subsection{Skill 飞轮：越用越强}

Skill 是给 agent 复用的横向能力（playbook），目前内置 65 个。能力缺失时由 distiller 在线蒸馏出新 skill 并回写。\textbf{每一次失败的研究尝试，都会沉淀出可复用的 skill、闸门、benchmark 代码与审计产物}，让下一次自主研究更强——这是一个随使用量自我增强的数据/能力飞轮。

%======================================================================
% 壁垒
%======================================================================
\section{竞争壁垒（Moat）}

\begin{tabularx}{\textwidth}{@{}>{\raggedright\arraybackslash}p{4.2cm} X@{}}
\toprule
\headrow \thd{壁垒} & \thd{为什么难被复制} \\
\midrule
\altrow 长程记忆引擎 & 反遗忘上下文治理是数十个版本踩坑迭代出来的工程 know-how，不是调个 prompt 能复刻的。 \\
诚信护栏体系 & 一整套防造假、可审计、强复现的硬闸门，需要大量真实跑批中的失败案例反推。 \\
\altrow 反平庸评审 & Reviewer agent + 阶段 checklist 体系，是"完成度判定"的事实来源，竞品多用脆弱的关键词/正则。 \\
Skill 飞轮 & 越用越强的能力库，先发者积累的 skill / 闸门 / benchmark 形成数据护城河。 \\
\altrow 架构哲学 & "harness 不替 agent 判断"让系统能力随底层模型升级自动抬升，不被写死规则封顶。 \\
\bottomrule
\end{tabularx}

%======================================================================
% 市场
%======================================================================
\section{市场机会}

\subsection{市场分层（TAM / SAM / SOM）}

\begin{tabularx}{\textwidth}{@{}l >{\raggedright\arraybackslash}p{3.2cm} X@{}}
\toprule
\headrow \thd{层级} & \thd{规模（估算）} & \thd{口径} \\
\midrule
\altrow TAM & 数千亿美元 & 全球研发（R\&D）支出 > 2.5 万亿美元/年中，可被"自主研究"替代或增益的知识工作部分。 \\
SAM & 数百亿美元 & 学术机构 + 企业研发实验室 + 量化/生物医药等高研究密度行业的研究人力与工具预算。 \\
\altrow SOM（3年） & 数亿美元 & 率先采用 AI 自主研究的头部实验室、AI 公司、量化基金、研发型企业的早期付费市场。 \\
\bottomrule
\end{tabularx}

\footnotesize\color{argusgray}注：以上为基于公开行业口径的数量级估算，用于说明市场量级，非精确测算。\normalsize\color{black}

\subsection{目标客户}
\begin{itemize}
  \item \hl{AI / 算法实验室}：把研究员从格式、产物记账中解放，聚焦 idea 与评估标准；用 Argus 做大规模 idea 筛选与基线复现。
  \item \hl{量化基金}：自主搜索、回测、验证因子与策略，端到端可审计——与诚信护栏天然契合。
  \item \hl{生物医药 / 材料研发}：高研究密度、强复现要求的实验科学。
  \item \hl{高校与研究生培养}：把导师精力从"改格式"转向"判断 thesis 质量与评估标准"。
\end{itemize}

%======================================================================
% 商业模式
%======================================================================
\section{商业模式与定价}

Argus 采用\textbf{多层组合变现}，核心是把"研究产出"而非"软件席位"作为价值锚：

\begin{tabularx}{\textwidth}{@{}>{\raggedright\arraybackslash}p{3.6cm} X >{\raggedright\arraybackslash}p{3.4cm}@{}}
\toprule
\headrow \thd{模式} & \thd{说明} & \thd{定价锚（建议）} \\
\midrule
\altrow 按研究产出计费 & 每完成一篇通过评审认证的研究/论文计费 & 单篇 US\$2,000–5,000 \\
团队订阅（SaaS） & 按"研究席位"包月，含一定额度的自主研究运行 & US\$1,500–3,000 /席位/月 \\
\altrow 私有化部署 & 面向量化/药企/大厂的本地部署 + 定制领域 skill & 年度许可 + 实施费 \\
平台 / API & 长程 agent 记忆与编排能力对外开放 & 用量计费 \\
\bottomrule
\end{tabularx}

\textbf{成本结构由真实数据支撑（非拍脑袋）。} 基于仓库内权威成本口径：\texttt{gpt-5.4} 输入 US\$1.25/M、命中缓存输入 US\$0.125/M、输出 US\$10/M。在重复跑批场景下，Engineer 的 prompt cache 实测命中率约 \hl{95\%}（单次最大成本杠杆），使长任务的 token 成本下降约 74\%。系统内置单任务预算上限（默认 US\$30）与每日上限（默认 US\$180）做硬性成本护栏。

%======================================================================
% 单位经济
%======================================================================
\section{单位经济模型（Unit Economics）}

以"一篇通过认证的研究/论文"为单位（数量级估算，随 benchmark 重度浮动）：

\begin{tabularx}{\textwidth}{@{}X r@{}}
\toprule
\headrow \thd{项目} & \thd{单篇估算（US\$）} \\
\midrule
\altrow API 成本（三层 agent，含 95\% 缓存杠杆） & 150 – 400 \\
真实实验 GPU 算力 & 50 – 300 \\
\altrow \textbf{单篇边际成本（COGS）} & \textbf{\textasciitilde 200 – 700} \\
\midrule
建议售价（按产出计费） & \textbf{2,000 – 5,000} \\
\altrow \textbf{毛利率} & \textbf{\textasciitilde 75\% – 90\%} \\
\bottomrule
\end{tabularx}

\begin{keybox}
\textbf{对照人力：} 一名研究生/RA 产出一篇可投稿论文，通常耗时数月、人力成本 US\$10,000–50,000。Argus 把同等产出的边际成本压到数百美元——\hl{10–100 倍的成本结构优势}，且可大规模并行。
\end{keybox}

%======================================================================
% 竞争格局
%======================================================================
\section{竞争格局}

\begin{tabularx}{\textwidth}{@{}>{\raggedright\arraybackslash}p{3.4cm} c c c c@{}}
\toprule
\headrow \thd{方案} & \thd{文献} & \thd{真实实验} & \thd{论文成稿} & \thd{诚信/复现闸门} \\
\midrule
\altrow 人类 RA / 研究生 & \faCheck & \faCheck & \faCheck & 依赖个人 \\
文献工具（Elicit 等） & \faCheck & \faTimes & \faTimes & \faTimes \\
\altrow 通用 Coding Agent & 部分 & 部分 & \faTimes & \faTimes \\
学术 AI Scientist 原型 & \faCheck & 玩具级 & 前70\% & 弱 \\
\rowcolor{arguslight}\textbf{Argus} & \faCheck & \textbf{\faCheck 全量} & \textbf{\faCheck 可投稿} & \textbf{\faCheck 硬闸门} \\
\bottomrule
\end{tabularx}

\textbf{我们的差异化不在"能不能生成"，而在"能不能完成"——} 端到端跑通、真实证据、硬性复现，加上长程记忆引擎让它真的能撑过数百步而不退化。

%======================================================================
% GTM
%======================================================================
\section{进入市场策略（Go-to-Market）}

\begin{enumerate}
  \item \hl{灯塔案例（0→1）：} 先产出 1 篇端到端通过 Reviewer 全流程认证的旗舰论文，作为可信度证明（首要里程碑）。
  \item \hl{种子社区（1→10）：} 定向邀请头部 AI 实验室 / 高校课题组试用，建立 mentor-in-the-loop 检查点（idea 短名单、实验设计、结论校准、终稿 go/no-go）。
  \item \hl{行业纵深（10→100）：} 切入量化金融与生物医药——这两个行业对"可审计、可复现"付费意愿最强，与诚信护栏天然契合。
  \item \hl{平台化：} 把长程记忆与编排能力作为"长程 agent 平台"对外开放，超越科研单一场景。
\end{enumerate}

%======================================================================
% 路线图
%======================================================================
\section{路线图与里程碑（融资后 18 个月）}

\begin{tabularx}{\textwidth}{@{}l >{\raggedright\arraybackslash}p{3.0cm} X@{}}
\toprule
\headrow \thd{时点} & \thd{里程碑} & \thd{关键交付} \\
\midrule
\altrow M+3 & 灯塔论文 & 首篇 research→submission 端到端 Reviewer 认证论文；上线 daemon 状态/证据缺口实时看板。 \\
M+6 & 种子客户 & 10 家实验室/团队试用；mentor-in-the-loop 检查点；扩展 benchmark 来源至公开前沿套件。 \\
\altrow M+12 & 商业化启动 & 多领域可复用入口（NLP/CV/AI-infra）；首批付费客户；ARR 突破人民币 300 万元。 \\
M+18 & A 轮就绪 & ARR run-rate 人民币 1,500 万元；私有化部署样板客户；平台 API 内测。 \\
\bottomrule
\end{tabularx}

%======================================================================
% 团队
%======================================================================
\section{团队}

\textbf{（待补充——以下为结构占位，请按实际填写）}
\begin{itemize}
  \item \hl{创始人 / CEO：} \_\_\_\_\_\_\_\_——研究方向、过往成果、为何是做这件事的合适人选。
  \item \hl{技术负责人 / CTO：} \_\_\_\_\_\_\_\_——长程 agent / 大模型系统工程背景。
  \item \hl{核心研究/工程：} \_\_\_\_\_\_\_\_——已交付完整系统（153 测试、CI 全绿、65 内置 skill）的工程能力即最好背书。
  \item \hl{顾问：} \_\_\_\_\_\_\_\_——学术界（评审标准）与行业（量化/药企）顾问。
\end{itemize}

%======================================================================
% 财务
%======================================================================
\section{财务预测（三年 · 目标值）}

\begin{tabularx}{\textwidth}{@{}X r r r@{}}
\toprule
\headrow \thd{指标（人民币）} & \thd{2026 (Y1)} & \thd{2027 (Y2)} & \thd{2028 (Y3)} \\
\midrule
\altrow 营收 & 300 万 & 2,500 万 & 1.2 亿 \\
付费客户数 & 10 & 60 & 250+ \\
\altrow 毛利率 & \textasciitilde 70\% & \textasciitilde 78\% & \textasciitilde 82\% \\
团队规模 & 12 & 30 & 70 \\
\altrow 现金状态 & 烧钱（产品打磨） & 接近盈亏平衡 & 经营性现金流转正 \\
\bottomrule
\end{tabularx}

\footnotesize\color{argusgray}注：以上为基于当前产品成熟度与市场假设的\textbf{目标值/情景预测}，非承诺。实际取决于灯塔案例落地速度与付费转化。\normalsize\color{black}

%======================================================================
% The Ask
%======================================================================
\section{融资计划（The Ask）}

\begin{keybox}
\textbf{本轮：} 天使 / 种子轮，拟融资 \hl{人民币 1,500 万元}（约 US\$2.1M），出让 \hl{15\%–18\%}，对应投前估值约人民币 7,000 万–8,500 万元，可支撑 \hl{18 个月} runway。
\end{keybox}

\textbf{资金用途：}
\begin{tabularx}{\textwidth}{@{}X r@{}}
\toprule
\headrow \thd{用途} & \thd{占比} \\
\midrule
\altrow 研发团队（研究 + 长程 agent 工程） & 45\% \\
算力 / API（benchmark 跑批、模型推理） & 25\% \\
\altrow 商业化与 BD（种子客户、行业落地） & 20\% \\
合规、安全与运营 & 10\% \\
\bottomrule
\end{tabularx}

\textbf{投资人将获得：} 一个已建成、CI 全绿、具备稀缺长程记忆 know-how 的自主研究系统；一个随底层模型升级而自动增益、随使用而自我增强的能力飞轮；以及切入"AI 替代高水平研究人力"这一万亿级结构性机会的早期席位。

%======================================================================
% 风险
%======================================================================
\section{风险与应对}

\begin{tabularx}{\textwidth}{@{}>{\raggedright\arraybackslash}p{4.6cm} X@{}}
\toprule
\headrow \thd{风险} & \thd{应对} \\
\midrule
\altrow 学术伦理 / AI 论文接受度 & 定位"研究加速器 + mentor-in-the-loop"，全链路可审计、可溯源，人类保留终审；面向企业研发场景规避会议政策争议。 \\
对单一模型供应商依赖 & harness 已抽象多后端（backend 可切换），降低对单一供应商锁定。 \\
\altrow 幻觉 / 造假 & 诚信护栏是产品内核而非附加：真实 benchmark、禁止灌水、强制审计包。 \\
长程任务退化 / 失忆 & 自研反遗忘记忆引擎 + 结构化 session 滚动，已在数十版本迭代中验证。 \\
\altrow 成本失控 & 单任务/每日预算硬上限 + prompt cache 95\% 命中杠杆，成本可控可观测。 \\
\bottomrule
\end{tabularx}

\vspace{1.2em}
{\color{argusline}\rule{\textwidth}{1pt}}\\[0.4em]
{\sffamily\color{argusgray}\small 本文件包含前瞻性陈述与基于公开口径的估算，仅供本轮融资评估之用，不构成投资承诺或要约。市场规模与财务预测为情景假设，实际结果可能存在重大差异。\hfill 机密 · Confidential}
\end{document}
